\documentclass{ibetest}
\begin{document}
\maketest{Progress Test 1.C --- Elec-S1}
         {1.C \,$\cdot$\, Circuits and two-terminal components}
         {7 minutes}{14}

% ---------------------------------------------------------------------------
\exercise{Unit conversion}{1 mark}
Convert the power $P = \qty{2.5e-2}{\kilo\watt}$ into milliwatts. Give the exact result,
without rounding.
\answer{1}{2.2cm}{%
  $P = 2.5\times10^{-2}\times10^{3}\ \mathrm{W} = 25\ \mathrm{W}
      = 25\times10^{3}\ \mathrm{mW}$ \qquad
  $\boxed{P = 2.5\times10^{4}\ \mathrm{mW}}$}
\begin{scheme}
\pt{1}{1 pt}{$2.5\times10^{4}\ \mathrm{mW}$ (or $25\,000\ \mathrm{mW}$), unit written
  correctly.}
\end{scheme}

% ---------------------------------------------------------------------------
\exercise{Electrical power --- definition}{5 marks}
A two-terminal component, studied in the receiver convention, has a voltage $u(t)$ across
it and carries a current $i(t)$. Draw the receiver convention and give the power $p(t)$
received by the component, with its SI unit expressed in terms of volts and amperes. Then
give the energy it receives between $t_1$ and $t_2$, and its expression in the DC regime
($u = u_0$ and $i = i_0$ constant during $\Delta t$).
\answer{2,2,1}{6.4cm}{%
  \begin{minipage}[c]{0.26\linewidth}\centering\convfig{receiver}{D}\end{minipage}%
  \begin{minipage}[c]{0.72\linewidth}
    Receiver convention: the arrows of $u$ and $i$ point in \textbf{opposite} directions.
    \\[4pt]
    Power received: $\boxed{p(t) = u(t)\,i(t)}$ \ ($p > 0$ for a receiver)\\[4pt]
    Unit: the \textbf{watt}, $1\ \mathrm{W} = 1\ \mathrm{V}\times1\ \mathrm{A}$.
  \end{minipage}\\[6pt]
  Energy received: $\displaystyle E = \int_{t_1}^{t_2} p(t)\,\dd t
    = \int_{t_1}^{t_2} u(t)\,i(t)\,\dd t$ \quad (in joules)\\[3pt]
  DC regime: $\boxed{E = u_0\, i_0\, \Delta t}$}
\begin{scheme}
\pt{1}{2 pts}{Receiver convention drawn, arrows opposite (1~pt); $p = u\,i$ (1~pt).}
\pt{2}{2 pts}{$E = \int p\,\dd t$ (1~pt); $E = u_0 i_0 \Delta t$ in the DC regime (1~pt).}
\pt{3}{1 pt}{Watt, $\mathrm{W} = \mathrm{V\,A}$.}
\end{scheme}

% ---------------------------------------------------------------------------
\exercise{Reading a circuit diagram}{4 marks}
\begin{center}
\begin{circuitikz}[line width=0.6pt]
  \draw (0,0) to[battery1, invert, l=$e$] (0,3);
  \draw (0,3) -- (2.2,3);
  \draw (2.2,3) to[lamp, l_=$L_1$] (2.2,0);
  \draw (2.2,3) to[lamp, l=$L_2$] (4.8,3);
  \draw (4.8,3) to[lamp, l_=$L_3$] (4.8,0);
  \draw (4.8,3) -- (6.6,3);
  \draw (6.6,3) to[lamp, l=$L_4$] (6.6,1.5) to[lamp, l=$L_5$] (6.6,0);
  \draw (0,0) -- (2.2,0) -- (4.8,0) -- (6.6,0);
  \foreach \p in {(0,0),(2.2,0),(4.8,0),(6.6,0),(0,3),(2.2,3),(4.8,3),(6.6,3),(6.6,1.5)}
    \node[circ] at \p {};
\end{circuitikz}
\end{center}
Base your answers on the diagram exactly as drawn, even though the same circuit could be
drawn with fewer wires.
\question{How many nodes?}
\opts{0}{2}{1}{3}{0}{4}{0}{9}
\why{a node is a junction of at least three conductors, and all the points joined by wires
  alone are one and the same node (same potential). Distinct nodes: the top-left wire
  ($e$, $L_1$, $L_2$), the top-right wire ($L_2$, $L_3$, $L_4$) and the whole bottom wire
  ($e$, $L_1$, $L_3$, $L_5$). The dot between $L_4$ and $L_5$ joins only two conductors, so it
  is not a node. ``4'' counts the drawn junction dots: that was the trap in Test~1
  (expected answer 2, not 4).}
\question{How many branches?}
\opts{0}{4}{1}{5}{0}{6}{0}{11}
\why{a branch joins two nodes: $e$, $L_1$, $L_2$, $L_3$, and $L_4$ with $L_5$ in series
  (one single branch, since their common point is not a node).}
\question{How many two-terminal components, connecting wires included?}
\opts{0}{6}{0}{9}{1}{11}{0}{14}
\why{6 components ($e$, $L_1,\dots,L_5$) $+$ 5 connecting wires (each plain segment between
  two dots: 1 top-left, 1 top-right, 3 at the bottom) $= 11$. In Test~1 the answer was
  $4 + 6 = 10$ on the same principle.}
\question{How many independent meshes (independent mesh equations)?}
\opts{0}{2}{1}{3}{0}{4}{0}{5}
\why{$m = B - N + 1 = 5 - 3 + 1 = 3$: the three ``windows'' of the diagram.}

% ---------------------------------------------------------------------------
\exercise{Energy received by a phone}{4 marks}
A phone is charged in the DC regime: the charger applies $u = \qty{5.0}{V}$ and delivers
$i = \qty{2.0}{A}$ during $\Delta t = \qty{1}{h}\ \qty{30}{min}$. Compute the energy
received by the phone, in joules, then in watt-hours.
\data{$u = \qty{5.0}{V}$ \hspace{14mm} $i = \qty{2.0}{A}$ \hspace{14mm}
      $\Delta t = \qty{1}{h}\ \qty{30}{min}$}
\answer{1,1,1,1}{5.4cm}{%
  $E = u\,i\,\Delta t$ \qquad with \ $\Delta t = 1.5\ \mathrm{h} = 1.5\times3600\ \mathrm{s}
   = 5400\ \mathrm{s}$\\[3pt]
  $E = 5.0\ \mathrm{V}\times2.0\ \mathrm{A}\times5400\ \mathrm{s} = 10\ \mathrm{W}\times5400
   \ \mathrm{s}$ \qquad $\boxed{E = 5.4\times10^{4}\ \mathrm{J}}$\\[3pt]
  In watt-hours: $E = 10\ \mathrm{W}\times1.5\ \mathrm{h}$ \qquad
  $\boxed{E = 15\ \mathrm{W\,h}}$ \quad (check: $15\times3600 = 54\,000$)}
\begin{scheme}
\pt{1}{1 pt}{Literal expression $E = u\,i\,\Delta t$.}
\pt{2}{1 pt}{$\Delta t = 5400\ \mathrm{s}$.}
\pt{3}{1 pt}{$E = 5.4\times10^{4}\ \mathrm{J}$.}
\pt{4}{1 pt}{$E = 15\ \mathrm{W\,h}$.}
\end{scheme}
\begin{tip}
``1 h 30 min'' is $1.5\ \mathrm{h}$, not $1.3\ \mathrm{h}$. Convert times into seconds
before using SI formulas.
\end{tip}

\finish
\end{document}
